\documentclass[11pt, a4paper]{article}

\usepackage[utf8]{inputenc}
\usepackage[T1]{fontenc}
\usepackage{lmodern}
\usepackage{amsmath, amssymb, amsthm}
\usepackage[a4paper, margin=1in]{geometry}
\usepackage{microtype}
\usepackage[round]{natbib}
\usepackage{hyperref}
\usepackage{enumitem}
\usepackage{indentfirst}
\usepackage{xcolor}
\hypersetup{colorlinks=true, linkcolor=blue!50!black, urlcolor=blue!50!black, citecolor=blue!50!black}

\title{Wave Release Coordination under Vertical Resource Constraints in Multi-Story AMR Warehouses}
\author{Shiyue Hu \and (co-authors TBD)}
\date{May 2026}

\begin{document}

\maketitle

\section{Problem Formulation}\label{sec:problem}

We consider a multi-story warehouse where floor-bound AMRs share a small number of freight elevators across floors. Orders arrive over time, each with a known source floor and destination floor. A \emph{wave} is a subset of orders released into the system simultaneously at the start of a release window. The system performance metric is the wave makespan, the time at which the last order in the wave is delivered, given a fixed operational dispatch policy.

This setting decomposes naturally into two coupled decision layers. The \emph{tactical} layer chooses which orders to bundle into a wave $W$, namely the wave composition decision. The \emph{operational} layer determines, given $W$, which AMR picks which order and which AMRs ride which elevator together. We focus the analytical effort on the tactical layer and treat the operational layer through a fixed dispatch policy (capacity-bounded FIFO boarding; full simulator details appear in Section~5). The two layers are coupled through the makespan: a wave whose composition is structured to ease the downstream dispatch yields a smaller makespan than an unstructured wave of the same size.

\paragraph{Sets and indexes.}
\begin{itemize}[noitemsep, topsep=2pt, leftmargin=1.5em]
  \item $\mathcal{O}$: set of all orders, indexed by $o$
  \item $\mathcal{F} = \{1, 2, \ldots, F\}$: set of floors, indexed by $f$
  \item $\mathcal{A}$: AMR fleet, indexed by $a$
  \item $\mathcal{E}$: set of elevators, indexed by $e$
  \item $\mathcal{M} = \{M_1, M_2, M_3\}$: elevator-model set, indexed by $m$
\end{itemize}

\paragraph{Parameters.}
\begin{itemize}[noitemsep, topsep=2pt, leftmargin=1.5em]
  \item $F \in \mathbb{Z}_{\geq 2}$: number of floors
  \item $|\mathcal{A}|$: AMR fleet size
  \item $E$: number of elevators
  \item $c \in \mathbb{Z}_{\geq 1}$: per-trip elevator capacity
  \item $s_o \in \mathcal{F}$: source floor of order $o$
  \item $d_o \in \mathcal{F}, d_o \neq s_o$: destination floor of order $o$
  \item $r_o \in \mathbb{R}_{\geq 0}$: release time of order $o$
  \item $W_{\min}, W_{\max} \in \mathbb{Z}_{\geq 1}$: lower and upper bounds on wave cardinality
  \item $\tau_w \in \mathbb{R}_{\geq 0}$: start time of release window for wave $w$
  \item $\Delta \in \mathbb{R}_{> 0}$: length of release window
  \item $\sigma_{M_3} \in \{0.10, 0.20\}$: lognormal noise scale under stochastic model $M_3$
\end{itemize}

\paragraph{Decision variable.} The tactical decision variable is the wave composition $W \subseteq \mathcal{O}$, the subset of orders released together within the current wave window. An equivalent binary representation is $x_o \in \{0, 1\}$ with $x_o = 1$ iff $o \in W$.

\paragraph{Structured-feature representation.} Rather than working in the raw $|\mathcal{O}|$-dimensional decision space, we summarize each candidate wave by a three-dimensional fingerprint $\Phi(W) = (C(W), I(W), T(W))$:
\begin{itemize}[noitemsep, topsep=2pt, leftmargin=1.5em]
  \item \emph{Vertical concentration} $C(W) = 1 - H_{\text{dst}}(W) / \log_2 F$, where $H_{\text{dst}}(W)$ is the Shannon entropy of the destination-floor distribution within the wave; high $C$ means orders share few destination floors.
  \item \emph{Directional imbalance} $I(W) = \big| |W^{\uparrow}| - |W^{\downarrow}| \big| / |W|$, where $W^{\uparrow}$ and $W^{\downarrow}$ are the up- and down-going subsets; high $I$ means the wave is dominated by a single direction.
  \item \emph{Temporal clustering} $T(W) = \sigma_r(W) / \mu_r(W)$, the coefficient of variation of per-order release timestamps within the wave; high $T$ means orders arrive in tight bursts.
\end{itemize}
The three axes correspond to three orthogonal physical mechanisms (multi-floor spread, up-versus-down asymmetry, temporal bunching) by which wave composition affects vertical-resource stress. We treat $\Phi$ as a \emph{conceptual} decomposition rather than a predictive surrogate: each axis carries a clear physical unit and operational meaning.

\paragraph{Elevator models.} The dispatch policy is realized through a simulator under one of three elevator models:
\begin{itemize}[noitemsep, topsep=2pt, leftmargin=1.5em]
  \item $M_1$ (\emph{throughput aggregation}): each elevator is modeled as a server with a per-AMR throughput rate that absorbs batched-trip efficiency into a single multiplier.
  \item $M_2$ (\emph{true co-occupancy batching}): at each elevator, the next $c$ AMRs in the queue board together and ride as a unit, with realistic per-trip capacity, dwell time, and direction handling.
  \item $M_3$ (\emph{stochastic batching}): extends $M_2$ with lognormal noise on per-trip duration at scale $\sigma_{M_3}$, used in Section~5 for robustness sensitivity.
\end{itemize}
$M_1$ and $M_2$ are both well-attested in the warehouse OR literature; the \emph{structural} disagreement between throughput aggregation and explicit co-occupancy is the methodological gap our Hedge Rule (Section~4) resolves.

\paragraph{Objective function.} Minimize the expected wave makespan under the chosen elevator model:
\begin{equation}\label{eq:obj}\tag{1}
\min_{W \subseteq \mathcal{O}} \; \mathbb{E}\!\left[\mathcal{T}(W) \,\big|\, M\right],
\end{equation}
where $\mathcal{T}(W) := \max_{o \in W} t^{\text{deliver}}_o(W; M)$ is the wave makespan, $t^{\text{deliver}}_o(W; M)$ is the simulator-realized delivery time of order $o$ under model $M$, and the expectation is over operational stochasticity (AMR pickup times, elevator dwell times, and, under $M_3$, lognormal noise).

\paragraph{Constraints.} The wave composition $W$ is subject to the constraints (2)--(6) below: (2)--(5) are tactical-layer feasibility conditions on $W$, while (6) is operational feasibility realized through the simulator under model $M$ rather than as a formal optimization constraint.
\begin{gather}
W_{\min} \leq |W| \leq W_{\max}, \label{eq:c-card}\tag{2} \\
r_o \in [\tau_w, \tau_w + \Delta], \quad \forall o \in W, \label{eq:c-window}\tag{3} \\
\sum_{w} \mathbf{1}[o \in W_w] = 1, \quad \forall o \in \mathcal{O}, \label{eq:c-uniq}\tag{4} \\
x_o \in \{0, 1\}, \quad \forall o \in \mathcal{O}, \label{eq:c-dom}\tag{5} \\
n_{\text{aboard}}(e, t) \leq c, \quad \forall e \in \mathcal{E}, \; \forall t. \label{eq:c-op}\tag{6}
\end{gather}
Constraint (2) bounds wave cardinality between the warehouse's operating envelope: $W_{\min}$ prevents under-utilized release windows, and $W_{\max}$ respects picking-station throughput and elevator-queue saturation. Constraint (3) is temporal feasibility: each order's release time $r_o$ must lie inside the wave's release window of length $\Delta$. Constraint (4) requires the sequence of waves to partition $\mathcal{O}$, with no order skipped or double-served, where $W_w$ is the wave released at $\tau_w$ and $\mathbf{1}[\cdot]$ is the indicator function. Constraint (5) defines the binary inclusion indicator, equivalently $W = \{o \in \mathcal{O} : x_o = 1\}$. Constraint (6) caps the per-trip elevator load $n_{\text{aboard}}(e, t)$; the simulator additionally encodes three operational invariants not exposed as formal optimization constraints: single-order AMR carriage (one AMR carries at most one order at a time), source-before-destination floor sequencing (an AMR visits $s_o$ before $d_o$), and capacity-bounded FIFO queue discipline (AMRs board in arrival order, up to $c$). These invariants are realized through $\mathcal{T}(W)$'s simulator implementation.

Together, equation (1) subject to constraints (2)--(6) defines the \emph{wave release coordination problem under vertical resource constraints}. Direct solution of this stochastic optimization is challenging on two counts: $\mathcal{T}(W)$ is simulator-realized under $M$ rather than closed-form in $W$, and the decision space (subsets of $\mathcal{O}$) is combinatorial in $|\mathcal{O}|$. Section~4 therefore develops two analytical tools that operate on the structured-feature representation $\Phi(W)$ instead: a Bound-and-Gap decomposition theorem that bounds the value of structured information under any fixed feature partition of $\Phi$ space, and a Model-Dominance Hedge Rule that resolves structural disagreement among $\{M_1, M_2, M_3\}$ via a closed-form decision. The four problem-defining elements identified in Sections~1 and~2 (wave composition as decision variable, multi-story deployment, flexible AMR fleet, shared-elevator capacity coupling) are now jointly formalized within this two-stage scheduler.

\end{document}
